\section{Описание программы}

\subsection{Архитектура проекта}

Класс \texttt{ExamGrade} содержит два целочисленных поля: \texttt{student\_id} и \texttt{grade}. Конструктор по умолчанию обнуляет поля, конструктор с параметрами задаёт id и оценку. Равенство элементов проверяет оба поля. Оператор вывода печатает пару в круглых скобках.

Класс \texttt{ExamGradeSet} содержит указатель \texttt{data\_}, количество элементов \texttt{size\_} и вместимость \texttt{capacity\_}. Элементы хранятся в порядке добавления. Для поиска id применяется линейный проход. При первом добавлении выделяется место для одного элемента, далее вместимость удваивается. Удаление выполняется сдвигом последующих элементов влево.

\subsection{Управление памятью}

Копирующий конструктор выделяет отдельный массив и копирует элементы. Копирующее присваивание использует copy-and-swap. Перемещающие операции передают владение массивом, обнуляют указатель, размер и вместимость источника и объявлены \texttt{noexcept}. Перемещающее присваивание предварительно освобождает собственный массив и проверяет самоперемещение. Деструктор вызывает \texttt{delete[]}.

\subsection{Семантика операций}

\begin{itemize}
  \item \texttt{add} ищет id: новый элемент добавляется в конец, существующая оценка заменяется на переданную. При полном совпадении возвращается \texttt{false}.
  \item \texttt{contains}, \texttt{operator()} и \texttt{remove} работают по id, независимо от переданной оценки. Это выбранная трактовка операций для данного варианта.
  \item Объединение копирует левое множество и добавляет записи правого. При конфликте id используется оценка из правого множества.
  \item Пересечение сохраняет записи левого множества, id которых есть в правом; оценка берётся слева.
  \item Разность сохраняет записи левого множества, id которых отсутствуют в правом.
  \item Равенство множеств требует совпадения размеров, id и соответствующих оценок; порядок элементов не учитывается.
\end{itemize}

Политика выбора оценки при конфликте в объединении и пересечении в ТЗ не уточнена. Принятые правила обеспечивают однозначный результат. Все три операции создают новое множество и сохраняют исходные операнды.

Поиск, добавление и удаление имеют сложность $O(n)$. Перераспределение памяти при росте также требует $O(n)$. Для множеств размеров $n$ и $m$ пересечение и разность имеют сложность $O(nm)$, объединение --- $O(n+m(n+m))$. Сравнение множеств одинакового размера имеет сложность $O(n^2)$.

\subsection{Исходный код}

\subsubsection{Элемент множества}
\lstinputlisting[style=cppstyle,caption={src/exam\_grade.h}]{../../src/exam_grade.h}
\lstinputlisting[style=cppstyle,caption={src/exam\_grade.cpp}]{../../src/exam_grade.cpp}

\subsubsection{Класс множества}
\lstinputlisting[style=cppstyle,caption={src/exam\_grade\_set.h}]{../../src/exam_grade_set.h}
\lstinputlisting[style=cppstyle,caption={src/exam\_grade\_set.cpp}]{../../src/exam_grade_set.cpp}

\subsubsection{Точка входа}
\lstinputlisting[style=cppstyle,caption={src/main.cpp}]{../../src/main.cpp}

\subsubsection{Тесты}
\lstinputlisting[style=cppstyle,caption={tests/test\_exam\_grade\_set.cpp}]{../../tests/test_exam_grade_set.cpp}
